\documentclass[11pt,a4paper]{article}
\usepackage[T1]{fontenc}
\usepackage[utf8]{inputenc}
\usepackage[spanish,es-noshorthands]{babel}
\usepackage{amsmath,amssymb,bm}
\usepackage{graphicx,booktabs,hyperref,listings,geometry,enumitem}
\usepackage{siunitx}
\geometry{margin=2.4cm}
\graphicspath{{../../results/tmm/}{../../figures/}}
\hypersetup{colorlinks=true,linkcolor=black,citecolor=blue,urlcolor=blue}
\lstset{
  language=Python,
  basicstyle=\ttfamily\footnotesize,
  breaklines=true,
  columns=fullflexible,
  frame=single,
  framesep=3pt,
  showstringspaces=false,
  commentstyle=\itshape,
  keepspaces=true,
  inputencoding=utf8,
  extendedchars=true,
  literate={á}{{\'a}}1 {é}{{\'e}}1 {í}{{\'i}}1 {ó}{{\'o}}1 {ú}{{\'u}}1
           {ñ}{{\~n}}1 {Á}{{\'A}}1 {ó}{{\'o}}1 {μ}{{$\mu$}}1,
}
\lstdefinestyle{tiny}{basicstyle=\ttfamily\scriptsize,breaklines=true,columns=fullflexible}

\title{Solución numérica e implementación:\\
de las ecuaciones del PRB 2010 a las seis figuras}
\author{Emanuel Orozco Gallego}
\date{31 de agosto de 2026}

\setlength{\emergencystretch}{3em}
\begin{document}
\maketitle

\begin{abstract}
Este documento describe, con el código real del paquete \texttt{fibonacci\_photonics},
cómo se resuelve numéricamente el modelo de Reyes-Gómez \textit{et al.},
Phys.\ Rev.\ B \textbf{81}, 153101 (2010), y cómo se generan las Figs.~1--6.
La física de fondo está en la \textit{Guía conceptual} hermana; aquí el hilo
es el algoritmo: unidades, matriz de transferencia, barrido en frecuencia,
detección de subbandas y trazado.
\end{abstract}

\tableofcontents
\newpage

\section{Flujo global}

El cálculo no integra Maxwell en el tiempo. Para cada terna
\((m,\theta,\nu)\) evalúa la semitraza \(R_m(\omega,\theta)\) y decide si
esa frecuencia es banda o gap.

\begin{enumerate}
\item Un YAML en \texttt{configs/figure\_0N.yaml} fija \(a\), \(b\),
\(\omega_e/2\pi\), \(\omega_m/2\pi\), polarización y el barrido.
\item \texttt{config.load\_figure} valida el YAML con Pydantic y convierte milímetros y GHz a SI.
\item \texttt{TMMSolver(spec).scan} recorre \(\nu\), calcula \(R_m\) y
\(k L_m/\pi\).
\item \texttt{reproduction.run\_figure("figure\_0N", "tmm")} (o
\texttt{scripts/run\_figure.py}) dibuja y guarda PDF/PNG/SVG más los datos \texttt{.npz} y un \texttt{summary.json}.
\end{enumerate}

Dos caminos independientes para \(R_m\), que deben coincidir:

\begin{itemize}
\item Producto de matrices de capa: \(T_m=M_N\cdots M_1\),
\(R_m=\frac12\mathrm{Tr}(T_m)\).
\item Recurrencia de Kohmoto: \(R_0\), \(R_1\), \(R_2\) analíticos y
\(R_m=2R_{m-1}R_{m-2}-R_{m-3}\).
\end{itemize}

Por defecto las figuras usan \texttt{TMMSolver(spec, algorithm="recurrence")} (más barato en \(m\)
alto). El producto sirve de test.

\section{Arquitectura del código}

El paquete vive en \texttt{src/fibonacci\_photonics/}; la TMM está en
\texttt{solvers/tmm/} y la física compartida en \texttt{physics/}.

{\small
\begin{center}
\begin{tabular}{llp{6.8cm}}
\toprule
Módulo & Rol numérico & Qué calcula \\
\midrule
\texttt{physics/constants.py} & \(c\), GHz, mm & SI \\
\texttt{physics/fibonacci.py} & \(S_m\), \(F_m\), \(L_m\) & palabras y conteos \\
\texttt{physics/materials.py} & Drude y aire & \(\varepsilon(\omega)\), \(\mu(\omega)\) \\
\texttt{config/} & YAML \(\to\) SI & esquemas Pydantic, \texttt{SuperlatticeSpec} \\
\texttt{physics/electromagnetics.py} & \(n\), \(q\), \(Q\), \(\chi\) & cinemática de la onda \\
\texttt{solvers/tmm/transfer\_matrix.py} & \(M\), \(T_m\), \(R_m\) & TMM \\
\texttt{solvers/tmm/solver.py} & \(\cos(kL)=R\) & \texttt{TMMSolver}: barrido y evaluador \\
\texttt{solvers/scan.py} & bandas, intervalos & \texttt{DispersionScan}, máscara \\
\texttt{analysis/plasmon\_modes.py} & recuento & subbandas en una ventana \\
\texttt{analysis/bandwidth.py} & \(\nu_{\min}(\theta),\nu_{\max}(\theta)\) & Fig.~6 \\
\texttt{viz/} & estilo PRB & cortes de curva y relleno \\
\texttt{io/provenance.py} & metadatos & reproducibilidad \\
\bottomrule
\end{tabular}
\end{center}
}

Los parámetros científicos no están quemados en los scripts: están en YAML.

\section{Unidades}

El núcleo trabaja en SI. Las figuras hablan GHz y mm.

\begin{lstlisting}[style=tiny]
# physics/constants.py
SPEED_OF_LIGHT = 299_792_458.0  # m/s   [IMPL; el paper no da c]
GHZ = 1.0e9
MM  = 1.0e-3
\end{lstlisting}

Conversiones en \texttt{physics/units.py}; la tolerancia de banda
\(10^{-12}\) está en \texttt{solvers/scan.py}:

\begin{equation}
\omega=2\pi\nu_{\mathrm{GHz}}\times 10^9,\qquad
a_{\mathrm{SI}}=a_{\mathrm{mm}}\times 10^{-3}.
\end{equation}

Ejemplo: \(\omega_e/2\pi=3\,\mathrm{GHz}\) \(\Rightarrow\)
\(\omega_e=1.884955592\times 10^{10}\,\mathrm{rad/s}\).
Usar \(\omega_e=3\times 10^9\) sin el \(2\pi\) desplazaría todas las figuras.

\texttt{units.omega\_from\_nu\_ghz} es la única puerta de \(\nu\) a
\(\omega\) en el barrido.

\section{De YAML a \texttt{SuperlatticeSpec}}

Cada figura tiene su YAML. El de la Fig.~2, abreviado:

\begin{lstlisting}[style=tiny]
layer_a_thickness_mm: 12.0
layer_b_thickness_mm: 12.0
omega_e_over_2pi_ghz: 3.0
omega_m_over_2pi_ghz: 3.0
polarization: TE
fibonacci_orders: [3, 4, 5, 6]
thetas_rad: [1.0471975511965976]   # pi/3
frequency_min_ghz: 2.0
frequency_max_ghz: 4.0
frequency_points: 12000
plasmon_window_ghz: [2.0, 2.999]
\end{lstlisting}

\texttt{SuperlatticeConfig.to\_spec} construye aire (\(\varepsilon=\mu=1\)) y el Drude
con \(\omega_e\), \(\omega_m\) ya en rad/s.
\texttt{load\_figure("figure\_02")} devuelve la configuración validada
(\texttt{ModeCountConfig}), el \texttt{spec} y la ruta: la configuración
conserva ángulos, ventanas y número de puntos que el \texttt{spec} no necesita.
Un valor inválido (espesor negativo, ventana invertida, \(\theta>\pi/2\),
clave desconocida) detiene la carga con un \texttt{ConfigError}.

\section{Fibonacci en código}

\texttt{paper\_fibonacci(m)} implementa \(F_0=F_1=1\), no \(F_0=0\).
\texttt{sequence(m)} concatena \(S_m=S_{m-1}S_{m-2}\) con memoización
(\texttt{lru\_cache}).
Una segunda rutina, \texttt{sequence\_by\_inflation}, aplica \(B\to A\),
\(A\to AB\); los tests exigen que ambas den la misma palabra.

Conteos:

\begin{lstlisting}[style=tiny]
n_layers_a(m) = paper_fibonacci(m - 1)   # F(m-1)
n_layers_b(m) = paper_fibonacci(m - 2)   # F(m-2) = N_modos
cell_length   = n_a * a + n_b * b
\end{lstlisting}

Para \(m=0,1\) se usa \(F_{-1}=0\), \(F_{-2}=1\).
\texttt{substitution\_identity} comprueba
\(S_m(A,B)=S_{m-k}[S_{k+1},S_k]\).

El producto de matrices recorre \texttt{sequence(m)} de izquierda a derecha
y premultiplica: si \(S_2=AB\), el bucle hace \(T\leftarrow M_A I\) y luego
\(T\leftarrow M_B T\), es decir \(T=M_B M_A=T_0 T_1\), como el paper.

\section{Materiales}

Aire: \(\varepsilon\) y \(\mu\) constantes.
Drude: se rechaza \(\omega=0\) (polo). No hay \(\gamma\).

\begin{lstlisting}[style=tiny]
eps = epsilon_0 - omega_e**2 / omega**2
mu  = mu_0     - omega_m**2 / omega**2
\end{lstlisting}

\(\nu_e=\omega_e/(2\pi\sqrt{\varepsilon_0})\) y análogo para \(\nu_m\).
Con \(\varepsilon_0=\mu_0=1\) coinciden con \(\omega_{e,m}/2\pi\) en GHz.

\section{Cinemática: \(n\), \(q\), \(Q\), \(\chi\)}

Índice: \(n=\sqrt{\mu}\sqrt{\varepsilon}\) (dos raíces principales de NumPy).
Para \(Q\) se usa \(n^2=\mu\varepsilon\), que no necesita el signo de \(n\).

\begin{lstlisting}[style=tiny]
q = (omega / c) * n_A * sin(theta)          # omega/c restaurado
Q = sqrt( (omega/c)**2 * n**2 - q**2 )      # complejo
chi = mu   si TE;   chi = epsilon si TM
\end{lstlisting}

\texttt{sqrt} complejo cubre ondas propagantes (\(Q\) casi real) y evanescentes
(\(Q\) casi imaginario). \(\cos(Qd)\) y \(\sin(Qd)\) de NumPy aceptan argumento
complejo y reproducen \(\cosh/\sinh\) sin un \texttt{if}.

\section{Matriz de una capa}

En un medio homogéneo, \(E=A\cos(Qz)+B\sin(Qz)\). Exigir
\(\Psi(z+d)=M\Psi(z)\) con \(\Psi=(E,\,\chi^{-1}E')\) da

\begin{equation}
M=\begin{pmatrix}
\cos(Qd) & (\chi/Q)\sin(Qd)\\
-(Q/\chi)\sin(Qd) & \cos(Qd)
\end{pmatrix}.
\end{equation}

Cuando \(\lvert Q\rvert<\texttt{SMALL\_Q\_ABS}\), \(\sin(Qd)/Q\to d\) se
reemplaza por el desarrollo \(d(1-\delta^2/6)\) para no dividir por cero
(\(Q=0\) es el umbral de corte).

\texttt{layer\_matrix} vectoriza sobre el barrido de \(\omega\): el array
tiene forma \((2,2,N_\omega)\). El producto es
\texttt{einsum("ij...,jk...->ik...")}.

Comprobaciones analíticas (tests):

\begin{itemize}
\item \(\det M=1\).
\item \(\frac12\mathrm{Tr}(M_B)=R_0=\cos(Q_B b)\).
\item \(\frac12\mathrm{Tr}(M_A)=R_1=\cos(Q_A a)\).
\item \(\frac12\mathrm{Tr}(M_B M_A)=R_2\) (ecuación 9 del paper).
\end{itemize}

\section{Semitraza: producto y recurrencia}

\texttt{analytic\_r0\_r1\_r2} escribe (7)--(9) con \(\chi\) en lugar de \(\mu\),
para que TM sea el mismo código.

\texttt{semitrace\_by\_recurrence} itera
\begin{equation}
R \leftarrow 2 R_{m-1} R_{m-2}-R_{m-3}
\end{equation}
desde \(m=3\) hasta el orden pedido.

\texttt{semitrace\_by\_product} hace \(\frac12\mathrm{Tr}(T_m)\) con el
producto de capas. Los tests piden acuerdo relativo \(\sim 10^{-8}\) hasta
\(m=7\).

También existe \texttt{cell\_matrix\_by\_recurrence\_blocks}:
\(T_m=T_{m-2}T_{m-1}\) con \(T_0=M_B\), \(T_1=M_A\), para contrastar el
orden de multiplicación.

\section{Barrido de dispersión}
\label{sec:scan}

Función central: \texttt{TMMSolver.scan}. Antes de calcular valida que las
frecuencias sean finitas, positivas y estrictamente crecientes, que
\(\theta\in[0,\pi/2]\) y que \(m\ge 0\) sea entero.

\begin{enumerate}
\item \(\omega_i=2\pi\nu_i\times 10^9\) en una malla uniforme de \(\nu\).
\item Se marcan polos de Drude
\(\lvert\omega-\omega_e\rvert<10^{-6}\omega_e\) o igual con \(\omega_m\);
esos puntos se excluyen de las bandas (\(\Psi\) no está definido si
\(\mu=0\) o \(\varepsilon=0\)).
\item Se evalúa \(R_m(\omega_i)\) por recurrencia o producto.
\item Máscara de banda: \(\lvert R\rvert\le 1+10^{-12}\). El entero
\texttt{clipped\_near\_unit} cuenta cuántos puntos cayeron en
\((1,1+10^{-12}]\) (redondeo, no recorte silencioso del gap).
\item Donde hay banda,
\(k L_m/\pi=\arccos(\mathrm{Re}\,R)/\pi\in[0,1]\).
En gaps se deja \texttt{NaN}, no se inventa un \(k\).
\end{enumerate}

El objeto \texttt{DispersionScan} guarda \(\nu\), \(\omega\), \(R\),
\(\lvert R\rvert\), la máscara, \(k L_m/\pi\), el método y cuántos puntos se
cliparon.

\texttt{allowed\_intervals} recorre la máscara y agrupa componentes conexas:
cada racha de \texttt{True} es un \texttt{FrequencyInterval}
\((\nu_{\min},\nu_{\max},n_{\mathrm{puntos}})\), con
\(\Delta\nu=\nu_{\max}-\nu_{\min}\).

\section{Detección de modos plasmon-polaritón}

\texttt{detect\_plasmon\_modes} filtra los intervalos que intersectan una
ventana \([\nu_a,\nu_b]\) y tienen al menos \texttt{min\_points} muestras
(para no contar un punto aislado de malla como modo).
El número esperado es \(n\_layers\_b(m)=F_{m-2}\).

\texttt{merge\_touching\_intervals} une dos intervalos si el hueco entre ellos
es menor que \texttt{gap\_ghz} (artefacto de un punto prohibido en medio de
una banda).

En la Fig.~2 la ventana es \([2.0,\,2.999]\,\mathrm{GHz}\): se corta justo
debajo del polo \(\nu_m=3\,\mathrm{GHz}\).

\section{Cómo se dibuja \(k(\nu)\) sin manchas}

Un \texttt{plot(k, nu)} ingenuo une puntos consecutivos en \(\omega\).
Cuando \(k\) salta de \(\sim 0\) a \(\sim 1\) al cruzar un gap, Matplotlib
traza una recta horizontal que rellena el panel (el artefacto de la primera
Fig.~1).

\texttt{plot\_dispersion\_branches} hace tres cosas:

\begin{enumerate}
\item Pone \texttt{NaN} donde no hay banda.
\item Cierra huecos de \emph{un} solo punto si los vecinos en \(k\) están
cerca (\texttt{\_fill\_isolated\_nans}): quita un poro numérico en \(k=0\).
\item Inserta \texttt{NaN} si \(\lvert\Delta k\rvert>0.12\)
(\texttt{\_polyline\_with\_breaks}): corta saltos de zona de Brillouin.
\item Dibuja \(+k\) y \(-k\) para llenar \([-1,1]\).
\end{enumerate}

Estilos: \(S_3\) rojo discontinuo, \(S_4\) azul continuo, ticks hacia dentro.

\section{Pipeline de cada figura}

\subsection{Figura 1}

Script: \texttt{scripts/tmm/figure\_01.py}.
Config: \(\nu_e=\nu_m=3\,\mathrm{GHz}\), \(m\in\{3,4\}\),
\(\theta\in\{0,\pi/12,\pi/6,\pi/3\}\),
\(12000\) puntos en \([0.15,5]\,\mathrm{GHz}\) (se evita \(\nu=0\), polo de
Drude).

Para cada panel: \texttt{TMMSolver.scan} con la recurrencia,
\texttt{plot\_dispersion\_branches} con \texttt{S3\_STYLE}/\texttt{S4\_STYLE},
ejes \(k L_m/\pi\in[-1,1]\), \(\nu\in[0,5]\).
Se guarda también la frecuencia donde se cierra el gap \(\langle n\rangle=0\):
\texttt{zero\_average\_index\_frequency\_ghz}.

\subsection{Figura 2}

\(\theta=\pi/3\), \(\nu\in[2,4]\,\mathrm{GHz}\), \(m=3,4,5,6\).
Después de dibujar se llama a \texttt{detect\_plasmon\_modes} en
\([2,2.999]\,\mathrm{GHz}\).
El \texttt{summary.json} debe leer \texttt{detected = expected = 1,2,3,5}.
Línea horizontal a \(\nu_m=3\,\mathrm{GHz}\).

\subsection{Figura 3}

Igual que la Fig.~1 en geometría y ángulos, pero
\(\omega_m/2\pi=1\,\mathrm{GHz}\).
No se dibuja línea de referencia (el original no la tiene).
La banda plana a \(1\,\mathrm{GHz}\) aparece sola, como un segmento casi
horizontal, cuando \(\theta\neq 0\).

\subsection{Figura 4}

Cuatro ventanas densas (\(16000\) puntos cada una):

\begin{center}
\begin{tabular}{lccc}
\toprule
Panel & \(m\) & \(\theta\) & \(\nu\) (GHz) \\
\midrule
(a) & 3 & \(\pi/12\) & \(0.994\)--\(1.002\) \\
(b) & 3 & \(\pi/3\) & \(0.94\)--\(1.02\) \\
(c) & 4 & \(\pi/12\) & \(0.994\)--\(1.002\) \\
(d) & 4 & \(\pi/3\) & \(0.94\)--\(1.02\) \\
\bottomrule
\end{tabular}
\end{center}

Se detectan intervalos y se informa \(\Delta\nu\).
Valores de esta implementación: \(4.51\,\mathrm{MHz}\) (\(m=3\), \(\pi/12\))
y \(33.1\,\mathrm{MHz}\) (\(m=3\), \(\pi/3\)).
Convergencia (\texttt{scripts/tmm/convergence.py}): el ancho a
\(\pi/3\), \(m=3\), pasa de \(0.03299\,\mathrm{GHz}\) (500 puntos) a
\(0.03310\,\mathrm{GHz}\) (16000); estable desde 2000 puntos (\(<0.2\%\)).

\subsection{Figura 5}

Para \(\theta=\pi/12\) y \(\pi/3\), y cada \(m=2,\ldots,8\):
escanear, extraer intervalos en la ventana, dibujar un palo vertical
\([m,m]\times[\nu_{\min},\nu_{\max}]\) con
\texttt{BRANCH\_COLORS[índice]} (negro, rojo, verde, azul, \ldots de menor a
mayor \(\nu\)).
Eje \(x\): ``Fibonacci order''. Eje \(y\): ``bandwidth (GHz)'', rotulado como
el impreso; los valores son frecuencias.
Aquí no se usa \texttt{merge\_touching\_intervals}: a \(\theta=\pi/12\) y
\(m=7\) hay gaps físicos (\(|R_m|>1\)) de apenas \(0.4\,\mathrm{kHz}\), y
fusionarlos reduciría el conteo por debajo de \(F_{m-2}\).
Con la detección directa se obtienen \(1,1,2,3,5,8,13\) subbandas para
\(m=2,\ldots,8\) en ambos ángulos, estable con 20\,000, 80\,000 y 320\,000
puntos.

\subsection{Figura 6}

No es un \(\Delta\nu(\theta)\) en líneas. Algoritmo:

\begin{enumerate}
\item Malla de \(\theta\in[0,\pi/3]\) (41 puntos). La malla en \(\nu\) une
60\,000 puntos logarítmicos en \(\nu_m-\nu\in[10^{-10},\,0.1]\,\mathrm{GHz}\)
y 250\,000 puntos uniformes en \([0.95,\,\nu_m)\) (paso de
\(0.2\,\mathrm{kHz}\)).
Hace falta porque cerca de \(\theta=0\) las subbandas quedan a
\(\sim10^{-5}\,\mathrm{GHz}\) de \(\nu_m\) con anchos de apenas
\(\sim5\times10^{-8}\,\mathrm{GHz}\), y hacia \(\theta\approx\pi/12\) hay gaps físicos
de \(0.4\,\mathrm{kHz}\).
Una malla uniforme de 16\,000 puntos cortada en \(\nu_m-10^{-4}\) perdía todas
las subbandas por debajo de \(\sim3^\circ\) y parte de ellas hasta
\(\sim20^\circ\).
\item \texttt{bandwidth\_versus\_angle} sin fusionar intervalos: en cada
\(\theta>0\) se obtienen exactamente \(F_{m-2}\) subbandas conexas (el
script avisa si algún ángulo no las tiene).
\item Cada borde se refina con Brent entre la última muestra con
\(|R_m|\le1\) y su vecina exterior, hasta donde \(|R_m|=1\)
(\texttt{refine=True}). Sin este paso el ancho queda subestimado hasta en dos
pasos de malla, lo que importa en las subbandas más estrechas: con \(m=7\)
algunas perdían más del 1\,\% de su ancho. Lo mismo se hace con los bordes de
la Fig.~5.
\item Como el número de subbandas no cambia con \(\theta\), la subbanda
\(i\)-ésima contada desde abajo es la misma rama física en todos los ángulos
y conserva su color.
\item A \(\theta=0\) el plasmón está en el gap: todas las ramas se clavan en
\((\theta,\nu)=(0,\nu_m)\), como en el original.
\item \texttt{fill\_between} por tramos contiguos (no se unen huecos con una
diagonal).
\end{enumerate}

Colores, de menor a mayor \(\nu\): negro, rojo, verde, azul, magenta, cian.
Número de regiones para todo \(\theta>0\): \(1,1,2,3,5,8=F_{m-2}\) para
\(m=2,\ldots,7\).

\section{Regularizaciones numéricas}

{\small
\begin{center}
\begin{tabular}{llp{5.8cm}}
\toprule
Situación & Qué hace el código & Por qué \\
\midrule
\(\omega=0\) & \texttt{ZeroDivisionError} & Drude diverge \\
\(\omega=\omega_e\) o \(\omega_m\) & vecindad relativa \(10^{-6}\) excluida & polo de \(\varepsilon\) o \(\mu\) \\
\(Q\to 0\) & \(\sin(Qd)/Q\to d(1-\delta^2/6)\) & corte geométrico \\
\(\lvert R\rvert=1+\varepsilon_{\mathrm{maq}}\) & banda si \(\varepsilon\le 10^{-12}\) & redondeo, documentado \\
\(\mathrm{Im}\,R\neq 0\) & se usa \(\mathrm{Re}\,R\) en \(\arccos\) & modelo sin pérdidas \\
\bottomrule
\end{tabular}
\end{center}
}

A \(\theta=0\) con \(\omega_e=\omega_m\), \(\lvert R_m\rvert\le 1\) en todo el
barrido: no es un bug, es adaptación de impedancia con el aire.

\section{Tests automáticos}

\texttt{pytest tests} cubre:

\begin{itemize}
\item \texttt{test\_fibonacci.py}: palabras, inflación, \(N_B\), \(L_m\),
identidad de substitución.
\item \texttt{test\_materials.py}: ceros de Drude, \(\nu_e\) y \(\nu_m\).
\item \texttt{test\_transfer\_matrix.py}: \(\det T=1\), (7)--(9), recurrencia
vs producto, y que TE coincide con TM a \(\theta=0\).
\item \texttt{test\_dispersion.py}: máscara de banda, frecuencia
\(\langle n\rangle=0\).
\item \texttt{test\_validation.py}: \(N_{\mathrm{modos}}=F_{m-2}\) en las
ventanas de las Figs.~2 y 4.
\item \texttt{test\_params.py}: YAML de cada figura.
\end{itemize}

Tres casos se omiten (\texttt{skip}) cuando \(k>m\) en una identidad de
Fibonacci que no aplica.

\section{Cómo reproducir las gráficas}

\begin{lstlisting}
python3 -m venv .venv
source .venv/bin/activate
pip install -r requirements.txt
pytest
python scripts/tmm/run_all.py
# o una sola: python scripts/run_figure.py 4 --method tmm
\end{lstlisting}

Salida de cada figura:

\begin{itemize}
\item \texttt{results/tmm/figure\_0N/output/figure\_0N.pdf} (y PNG, SVG).
\item \texttt{results/tmm/figure\_0N/data/*.npz}: \(\nu\), \(\lvert R\rvert\),
máscara, \(k L_m/\pi\).
\item \texttt{results/tmm/figure\_0N/data/summary.json}: recuentos y anchos.
\end{itemize}

\texttt{make figures} (o \texttt{make tmm}) lanza además
\texttt{scripts/tmm/convergence.py}. La verificación con ondas planas
(\texttt{scripts/pwe/}, \texttt{scripts/comparison/}) está descrita en el
documento \textit{El método de expansión en ondas planas} y en el anexo
\textit{Comparación TMM vs PWE}.

Para TM: \texttt{polarization: TM} en el YAML (no hay figuras TM en el
original; el código está listo).
Para otro plasma: editar \texttt{omega\_e\_over\_2pi\_ghz} y
\texttt{omega\_m\_over\_2pi\_ghz}.
Para otro orden: \texttt{fibonacci\_orders}.

\section{Decisiones que no están en el paper}

Todas etiquetadas en comentarios como implementación, no como dato original:

\begin{center}
\begin{tabular}{ll}
\toprule
Decisión & Valor \\
\midrule
\(q\) & \((\omega/c)n_A\sin\theta\) \\
\(c\) & \(299792458\,\mathrm{m/s}\) \\
matriz de capa & \(M(Q,d,\chi)\) unimodular \\
rama de \(Q\) & \texttt{sqrt} complejo principal \\
banda & \(\lvert R\rvert\le 1+10^{-12}\) \\
ancho & \(\nu_{\max}-\nu_{\min}\) de cada componente conexa \\
Fig.~6 & parámetros de Figs.~3--5, \(m=2\ldots 7\) \\
\bottomrule
\end{tabular}
\end{center}

No se calibró nada contra el dibujo impreso.

\section{Mapa de scripts a funciones}

\begin{center}
\begin{tabular}{lll}
\toprule
Figura & Script & Núcleo \\
\midrule
1 & \texttt{tmm/figure\_01.py} & \texttt{TMMSolver.scan} + \texttt{plot\_dispersion\_branches} \\
2 & \texttt{tmm/figure\_02.py} & idem + \texttt{detect\_plasmon\_modes} \\
3 & \texttt{tmm/figure\_03.py} & igual que 1, otro YAML \\
4 & \texttt{tmm/figure\_04.py} & ventanas densas + anchos \\
5 & \texttt{tmm/figure\_05.py} & \texttt{allowed\_intervals} por \(m\) \\
6 & \texttt{tmm/figure\_06.py} & \texttt{bandwidth\_versus\_angle} + \texttt{fill\_between} \\
\bottomrule
\end{tabular}
\end{center}

El camino mínimo, para una sola curva, es:

\begin{lstlisting}[style=tiny]
import numpy as np
from fibonacci_photonics import TMMSolver, load_spec

spec = load_spec("figure_02.yaml")
nu = np.linspace(2.0, 4.0, 4000)
scan = TMMSolver(spec).scan(m=3, theta=np.pi/3, nu_ghz=nu,
                            polarization="TE")
# scan.allowed, scan.k_lm_over_pi, scan.nu_ghz
\end{lstlisting}

Eso es, en código, toda la física de \(\cos(k L_m)=R_m\).

\begin{thebibliography}{9}
\bibitem{Reyes2010}
E.~Reyes-Gómez \textit{et al.}, Phys.\ Rev.\ B \textbf{81}, 153101 (2010).
\bibitem{Reyes2009EPL}
E.~Reyes-Gómez \textit{et al.}, EPL \textbf{88}, 24002 (2009).
\bibitem{Kohmoto1983}
M.~Kohmoto, L.~P.~Kadanoff y C.~Tang, Phys.\ Rev.\ Lett.\ \textbf{50}, 1870 (1983).
\end{thebibliography}

\end{document}
